% ============================================================================
% WORKBOOK — Section 4.2: Simplifying Expressions by Combining Like Terms
% CCSS 7.EE.A.1
% Practice-focused reinforcement for the study guide topic
% ============================================================================
\section{Combining Like Terms}
\topicTitle{Combining Like Terms}

\begin{quickReview}{Combining Like Terms}
\textbf{Like terms} have the \emph{same} variable raised to the \emph{same} power. Combine them by adding their \textbf{coefficients}.

\begin{itemize}[leftmargin=*]
    \item $3x$ and $7x$ are like terms $\to$ $3x + 7x = 10x$.
    \item Constants ($4$ and $-9$) are always like terms $\to$ $4 + (-9) = -5$.
    \item $5y$ and $5y^{2}$ are \textbf{not} like terms (different powers).
\end{itemize}

\medskip
\textbf{Example:} $4x + 3 - 2x + 8 = (4x - 2x) + (3 + 8) = 2x + 11$.
\end{quickReview}

% ── Warm-Up ──────────────────────────────────────────────────────────────────
\begin{practiceBox}[funGreen]{\faSun~Warm-Up}

\practiceHeader[funGreen]{Identify the like terms in each expression.}

\begin{multicols}{2}
\prob $5x + 3 + 2x$ \answerBlank[3cm]
\answerExplain{$5x$ and $2x$}{Both terms have the variable $x$ raised to the first power, so $5x$ and $2x$ are like terms.}

\prob $7a - 4 + 3b + 9$ \answerBlank[3cm]
\answerExplain{$-4$ and $9$}{The constants $-4$ and $9$ are like terms. $7a$ and $3b$ have different variables, so they are not like terms.}

\prob $6y + 2y^{2} - y$ \answerBlank[3cm]
\answerExplain{$6y$ and $-y$}{Both have the variable $y$ to the first power. $2y^{2}$ has a different exponent, so it is not a like term with $6y$.}

\prob $-3m + 8 + m - 5$ \answerBlank[3cm]
\answerExplain{$-3m$ and $m$; $8$ and $-5$}{$-3m$ and $m$ share the variable $m$. The constants $8$ and $-5$ are also like terms with each other.}

\prob $4 + 2n + 7n$ \answerBlank[3cm]
\answerExplain{$2n$ and $7n$}{Both terms have the variable $n$ to the first power. The constant $4$ stands alone.}

\prob $10 - 3p + 6p - 1$ \answerBlank[3cm]
\answerExplain{$-3p$ and $6p$; $10$ and $-1$}{The variable terms $-3p$ and $6p$ are like terms. The constants $10$ and $-1$ are like terms.}
\end{multicols}
\end{practiceBox}

% ── Practice A: Basic Simplifying ────────────────────────────────────────────
\begin{practiceBox}{Simplify Each Expression}

\practiceHeader{Combine like terms to simplify.}

\begin{multicols}{2}
\prob $8x + 3x$ \answerBlank[3cm]
\answerExplain{$11x$}{Both terms are like terms with variable $x$. Add the coefficients: $8 + 3 = 11$. Result: $11x$.}

\prob $5a - 9a$ \answerBlank[3cm]
\answerExplain{$-4a$}{Combine coefficients: $5 - 9 = -4$. Result: $-4a$.}

\prob $7y + 4 - 3y$ \answerBlank[3cm]
\answerExplain{$4y + 4$}{Combine the $y$-terms: $7y - 3y = 4y$. The constant $4$ stays. Result: $4y + 4$.}

\prob $-2n + 6 + 5n - 1$ \answerBlank[3cm]
\answerExplain{$3n + 5$}{Variable terms: $-2n + 5n = 3n$. Constants: $6 - 1 = 5$. Result: $3n + 5$.}

\prob $9 - 4b + 2b - 3$ \answerBlank[3cm]
\answerExplain{$-2b + 6$}{Variable terms: $-4b + 2b = -2b$. Constants: $9 - 3 = 6$. Result: $-2b + 6$.}

\prob $6x + 2 - 6x + 8$ \answerBlank[3cm]
\answerExplain{$10$}{Variable terms: $6x - 6x = 0$. Constants: $2 + 8 = 10$. The variable terms cancel out, leaving just $10$.}
\end{multicols}
\end{practiceBox}

% ── Practice B: Fractions and Decimals ───────────────────────────────────────
\begin{practiceBox}[funTeal]{Simplify with Fractions and Decimals}

\prob $\frac{1}{3}x + \frac{2}{3}x - 5$ \answerBlank[3cm]
\answerExplain{$x - 5$}{Add the coefficients: $\frac{1}{3} + \frac{2}{3} = \frac{3}{3} = 1$. Result: $1x - 5 = x - 5$.}

\prob $0.4n + 3.2 - 0.6n$ \answerBlank[3cm]
\answerExplain{$-0.2n + 3.2$}{Combine the $n$-terms: $0.4 - 0.6 = -0.2$. Result: $-0.2n + 3.2$.}

\prob $\frac{3}{4}y - \frac{1}{2}y + 7$ \answerBlank[3cm]
\answerExplain{$\frac{1}{4}y + 7$}{Find a common denominator: $\frac{3}{4}y - \frac{2}{4}y = \frac{1}{4}y$. Result: $\frac{1}{4}y + 7$.}

\prob $1.5a - 2 + 0.5a + 4.5$ \answerBlank[3cm]
\answerExplain{$2a + 2.5$}{Variable terms: $1.5a + 0.5a = 2a$. Constants: $-2 + 4.5 = 2.5$. Result: $2a + 2.5$.}

\end{practiceBox}

% ── Practice C: True or False ────────────────────────────────────────────────
\begin{practiceBox}[funPurple]{True or False?}

\trueOrFalse{$4x + 4y$ can be simplified to $8xy$.}
\answerExplain{False}{$4x$ and $4y$ are not like terms because they have different variables. You cannot combine them. The expression $4x + 4y$ is already simplified.}

\trueOrFalse{$3a + 5 - a = 2a + 5$.}
\answerExplain{True}{Combine the $a$-terms: $3a - a = 2a$. The constant $5$ stays. Result: $2a + 5$.}

\trueOrFalse{$7m - 7m = 0$.}
\answerExplain{True}{The terms are identical but opposite. Combining: $7m - 7m = 0$. When like terms cancel completely, the result is $0$.}

\end{practiceBox}

% ── Practice D: Error Analysis ───────────────────────────────────────────────
\begin{errorBox}{Spot the Mistake}

\prob Sam simplified $5x + 3 - 2x + 7$ and got $10x$. What went wrong?
\answerBlank[6cm]
\answerExplain{Sam combined unlike terms}{Sam incorrectly combined $5x$, $3$, $-2x$, and $7$ all together. The variable terms give $5x - 2x = 3x$ and the constants give $3 + 7 = 10$. The correct answer is $3x + 10$.}

\prob Mia says $3y + 2y^{2} = 5y^{3}$. Is she right? Explain.
\answerBlank[6cm]
\answerExplain{No — they are not like terms}{$3y$ has exponent $1$ and $2y^{2}$ has exponent $2$. Different exponents mean they are not like terms. Also, combining like terms adds coefficients, not exponents. The expression $3y + 2y^{2}$ is already simplified.}

\end{practiceBox}

% ── Challenge ────────────────────────────────────────────────────────────────
\begin{challengeBox}
\prob Simplify $3x + 2y - 5 + 4x - 7y + 12 - x$. \answerBlank[4cm]
\answerExplain{$6x - 5y + 7$}{Group by variable: $x$-terms: $3x + 4x - x = 6x$. $y$-terms: $2y - 7y = -5y$. Constants: $-5 + 12 = 7$. Result: $6x - 5y + 7$.}

\prob If $2x + 3y + 5 = 2x + 3y + k$, what is $k$?
\answerBlank[3cm]
\answerExplain{$5$}{Since $2x + 3y$ appears on both sides, they cancel out. This leaves $5 = k$, so $k = 5$.}
\end{challengeBox}

\encouragement{Great job simplifying expressions! Combining like terms is a skill you will use in every algebra topic ahead!}
